\section{Introduction and conventions}\label{sec:intro}

\subsection{The results and the scope of this paper}
A \emph{generalized Wigner matrix} is an $N\times N$ Hermitian or real symmetric random matrix $H=(h_{ij})$ whose entries are independent up to the symmetry
constraint, centred, with variances $\sigma_{ij}^2=\EE|h_{ij}|^2$ that may depend on $(i,j)$ but sum to one in every row. The paper
of Erd\H{o}s, Yau and Yin, \emph{Rigidity of eigenvalues of generalized Wigner matrices} (below: \emph{the original}), established four results for such
matrices, under a spectral gap condition on the matrix of variances and a subexponential decay condition on the entries:
\begin{itemize}[leftmargin=2em,itemsep=1pt]
\item the \emph{strong local semicircle law} (Theorem~\ref{thm:21}): the Stieltjes transform $m(z)$ of the empirical spectral measure and the individual entries $G_{ij}(z)$
of the resolvent are close to their deterministic limits, down to the optimal scale $\im z\gg N^{-1}$, with the optimal error $(N\im z)^{-1}$ for $m$, uniformly in
the energy up to the spectral edges, together with the bound $2+N^{-2/3+o(1)}$ on the norm of $H$;
\item the \emph{rigidity of eigenvalues} (Theorem~\ref{thm:22}): the $j$-th eigenvalue is within $N^{-2/3+o(1)}\min(j,N+1-j)^{-1/3}$ of its classical location $\gamma_j$;
\item the \emph{strong local ergodicity of Dyson Brownian motion} (Theorem~\ref{thm:23}): after a time $t\gg N^{-1}$ the local eigenvalue statistics in the bulk, averaged over a small
energy window, coincide with those of the Gaussian ensemble, with an explicit power-law error;
\item the \emph{universality of the extreme eigenvalues} (Theorem~\ref{thm:24}): the distributions of the largest (and of the smallest) eigenvalue of two generalized Wigner
matrices with matching first two moments agree on the scale $N^{-2/3}$.
\end{itemize}
The present paper, called NP below, proves these four theorems again. Its purpose is different from that of the original: NP is written as the mathematical
source of a formalization in Lean~4 with Mathlib. This has three consequences.

\emph{(1) Self-containedness.} No result is quoted from another paper or book. Whatever is not available in Mathlib is proved here. A statement marked \infra\ (``infrastructure'') is a standard result of analysis, linear algebra or probability
which is \emph{not} available in Mathlib and which NP therefore proves in full (for instance the Helffer--Sj\"ostrand formula, the joint eigenvalue density of the Gaussian ensembles,
the Pr\'ekopa--Leindler inequality); the marker carries no other meaning, and it indicates the places where the formalization will have to add material to Mathlib.
Appendix~\ref{app:mathlib} lists the facts that \emph{are} taken from Mathlib. In particular the inputs which the original takes from earlier papers
(the Schur complement identities, the large deviation bounds for quadratic forms, the properties of $\msc$, the a priori bound on the norm, the Helffer--Sj\"ostrand
argument for the counting function, the logarithmic Sobolev inequality and the local relaxation flow) are all proved, in most cases by an argument chosen for
its shortness and for the absence of analytic prerequisites.

\emph{(2) Finite-$N$ statements with explicit quantifiers.} There are no limits and no ``with high probability'': every probability is bounded by an explicit expression,
every constant has a stated dependence, and the order of the quantifiers is fixed once and for all (Convention~\ref{conv:const}).

\emph{(3) New proofs where the old ones are long.} The statements are those of the original, up to the corrections listed in Appendix~\ref{app:diff}; several proofs are new.

NP has three parts. Part~I (Sections~\ref{sec:intro}--\ref{sec:export}) contains the model, the four statements, and the proofs of Theorems~\ref{thm:21}
and~\ref{thm:22}. Part~II (Section~\ref{sec:edge}) proves Theorem~\ref{thm:24} and Part~III (Sections~\ref{sec:C1}--\ref{sec:C9}) proves Theorem~\ref{thm:23}; both use Theorems~\ref{thm:21}, \ref{thm:22} and the auxiliary results
collected in Section~\ref{sec:export}, and nothing else from Part~I. The appendices contain the bookkeeping of the exponents of Part~I (Appendix~\ref{app:ledger}), its dependency graph
(Appendix~\ref{app:dep}), the list of differences from the original (Appendix~\ref{app:diff}) and the list of facts taken from Mathlib (Appendix~\ref{app:mathlib}).

\subsection{Strategy of Part~I}
\begin{enumerate}[leftmargin=2em,itemsep=2pt]
\item \emph{One level parameter.} The original works with three coupled exponents (in the factors $(\log N)^{\ell}$, $(\log N)^{L}$ and in an iteration index).
Here a single real parameter $K$, the \emph{level}, with $(\log N)^3\le K\le N^{c_0}$, is fixed. Every error factor is an explicit power $K^{C}$ and every exceptional
probability is at most $N^{C}e^{-K}$. The main internal result is Theorem~\ref{thm:K}; Theorems~\ref{thm:21} and~\ref{thm:22}, with the original powers of $\log N$, follow
from it by the substitution $K=(\log N)^{\phi L}$ with a small constant $\phi$ (Proposition~\ref{prop:reduction}, proved in Section~\ref{sec:main}).
\item \emph{Probabilistic input} (Section~\ref{sec:prob}). Large deviation bounds for linear, quadratic and bilinear forms in independent variables with unequal, possibly
vanishing variances. They are derived from Hoeffding's inequality after truncation, with a deterministic decoupling of the off-diagonal quadratic form.
\item \emph{Entrywise (weak) law} (Section~\ref{sec:weak}). Schur complement identities, the large deviation bounds conditionally on a minor, stability of the self-consistent
equation (the vector part on $\be^\perp$ without any factor $\log N$; the scalar part is a quadratic inequality), and a \emph{deterministic} continuity argument in
$\eta=\im z$ along a finite lattice.
\item \emph{Fluctuation averaging} (Section~\ref{sec:FA}). For the index set $L$ of a $p$-fold product one conditions on the minor $H^{(L)}$; the block of $G$ on $L$ is the inverse of an
$|L|\times|L|$ matrix $M$ whose entry $M_{ab}$ depends only on the rows $a$ and $b$ of $H$. A finite Neumann expansion in the off-diagonal part of $M$ exhibits the exact dependence
structure, and the terms in which a lone index is not visited have zero expectation. This replaces the thirty pages of Section~7 of the original.
\item \emph{Strong law} (Section~\ref{sec:strong}). A moment inequality that improves itself in one step (no iteration), followed by a dichotomy and a continuity argument at the level
of events; a second pass gives the improved bound outside the spectrum.
\item \emph{Norm bound} (Section~\ref{sec:norm}). A crude bound $\norm H\le K^{C}$ from a grid net and Hoeffding's inequality; eigenvalues between $2+K^{C_3}N^{-2/3}$ and the crude bound
are excluded by the averaged law outside the spectrum.
\item \emph{Rigidity} (Sections~\ref{sec:det} and~\ref{sec:rig}). The Helffer--Sj\"ostrand formula, proved from the divergence theorem on rectangles (Lemma~\ref{lem:HS}); a deterministic
bound on the counting function in terms of $m-\msc$ (Lemma~\ref{lem:count}); and a deterministic conversion lemma which turns the bound on the counting function
into the bound on $\lambda_j-\gamma_j$ with the factor $\min(j,N+1-j)^{-1/3}$, at both edges at once (Lemma~\ref{lem:conv}).
\end{enumerate}
The semicircle law itself enters only through the algebraic function $\msc$: its equation, its size (Lemmas~\ref{lem:msc}, \ref{lem:mscsize}), and the fact that it is the
Stieltjes transform of $\rhosc$ (Lemma~\ref{lem:stieltjes}), which is needed only for the counting function.

\subsection{Conventions}\label{sec:conv}
\begin{convention}[Constants]\label{conv:const}
The \emph{model parameters} are $\vth>0$, $\delta_->0$, $\delta_+>0$, $C_0>0$. By the monotonicity of the assumptions
(Lemma~\ref{lem:mono}(ii)) we assume in all proofs that $\vth\le1$, $\dpm\le1$, $C_0\ge1$, and we put
\[
 \fa:=1/\vth\ \ (\ge1),\qquad \delta_{\min}:=\min(\delta_-,\delta_+).
\]
Letters $C,c$ denote positive constants that depend only on $(\vth,\delta_-,\delta_+,C_0)$; their values may change from
line to line, and they never depend on $N$, $K$, $z$, the ensemble, the field $\kk$, or indices. An \emph{absolute} constant is a fixed real number.
Named constants ($c_*$, $C_*$, $C_s$, $\lambda_{\rm st}$, $c_F$, $C_1,\dots,C_7$, \dots) keep their value throughout Part~I. Parts~II and~III use a few
constants with the same names ($C_2,C_3,C_4$ and $\Psi$ in Section~\ref{sec:edge}; $c_0,c_1$, $\Psi$ and others in Part~III); these are defined locally in the section where they occur,
are unrelated to the constants of Part~I, and this is recalled there. The letter $\delta$ without index is used only for the exponents of Theorems~\ref{thm:23} and~\ref{thm:24}
(as in the original); the two are unrelated, and each is local to its theorem and to its proof.
``For $N\ge N_0$'' means: for all $N$ above a threshold that depends only on the model parameters.
Quantifiers are always in the order: model parameters $\to$ constants $\to$ $N_0$ $\to$ ($N$, $\kk$, ensemble) $\to$ $K$ (or $L$) $\to$ $z$.
\end{convention}
\begin{convention}[Probability]\label{conv:prob}
All probabilities are explicit. An event is called \emph{$K$-likely} if its complement has probability at most $N^{C}e^{-K}$
with a constant $C$ as in Convention~\ref{conv:const}. A union of at most $N^{C'}$ complements of $K$-likely events has probability
at most $N^{C+C'}e^{-K}$, and
\[
 N^{C}e^{-K}\le e^{-K/2}\qquad\text{whenever }K\ge(\log N)^3\text{ and }(\log N)^2\ge2C,
\]
because then $C\log N\le\frac12(\log N)^3\le K/2$. Events indexed by uncountably many $z$
are always handled as: a finite lattice of $z$'s, a union bound, and a deterministic Lipschitz extension; the final statements
assert the existence of a measurable event of large probability on which a bound holds for all $z$.
\end{convention}
\begin{convention}[Spectral parameter]\label{conv:z}
$z=E+\ii\eta$ with $E\in\R$, $\eta>0$; $\kappa:=\big||E|-2\big|$; $\varkappa:=\sqrt{\kappa+\eta}$. We write $\C_+:=\{z\in\C:\im z>0\}$.
\end{convention}
\begin{convention}[General notation]\label{conv:not}
$\log$ is the natural logarithm. For $x\in\C^n$, $|x|$ is the Euclidean norm, $\norm x_\infty$ the maximum norm, and $\langle x,y\rangle=\sum_i\ol{x_i}y_i$;
for a matrix $A$, $\norm A$ is the operator norm with respect to the Euclidean norm. $\bone(\cdot)$ is an indicator, $\#S$ the cardinality of a finite set $S$,
$(t)_+:=\max(t,0)$, and $|I|$ is the length of an interval $I$. For a complex number $\zeta\notin(-\infty,0]$, $\sqrt\zeta$ is the principal square root (the one with positive real part).
Integrals without a measure are with respect to Lebesgue measure.
\end{convention}
